
The operators
    \texttt{\textbf{==}},
    \texttt{\textbf{!=}},
    \texttt{\textbf{is}},
    \texttt{\textbf{is not}},
    \texttt{\textbf{>}},
    \texttt{\textbf{>=}},
    \texttt{\textbf{<}},
    \texttt{\textbf{<=}}
    follow the 
    \href{https://docs.python.org/3.14}{\color{DarkBlue}Python 3.14 Documentation}.
    As in Python, \texttt{\textbf{True}} and \texttt{\textbf{False}} are
    represented by the integers 1 and 0, respectively, for the purpose of
    any comparisons. The operators \texttt{\textbf{is}} and \texttt{\textbf{is not}}
    test for identity. As in Python, we can assume that there exists only a single
    instance of each of \texttt{\textbf{None}}, \texttt{\textbf{True}},
    and \texttt{\textbf{False}} system-wide, and that all lists and functions have their
    own identity once created.

    The operators
    \texttt{\textbf{==}} and \texttt{\textbf{!=}} on lists compare corresponding
    elements $x$ and $y$ as \texttt{\textbf{$x$ is $y$ or $x$ == $y$}}, and so does
    \texttt{\textbf{in}}, as in Python: an element is equal to itself, even if it is
    not equal to itself on its own. In particular, a \texttt{\textbf{nan}} that is
    referred to twice is equal to itself inside a list:
    \texttt{[math\_nan] == [math\_nan]} evaluates to \texttt{\textbf{True}},
    whereas \texttt{math\_nan == math\_nan} evaluates to \texttt{\textbf{False}}.
    Whether two numbers that were computed separately are identical is left open, as in Python;
    for example, \texttt{\textbf{float("nan")}} twice gives two equal lists in PyPy, but not in CPython.

    List access \texttt{\textit{list}\texttt{\textbf{[}}\textit{index}\texttt{\textbf{]}}}
    and list assignment
    \texttt{\textit{list}\texttt{\textbf{[}}\textit{index}\texttt{\textbf{]} = }\textit{value}}
    require \textit{index} to evaluate to an \texttt{int}; as opposed to Python,
    \texttt{\textbf{True}} and \texttt{\textbf{False}} are not allowed as indices.
    The valid range of \textit{index} is \texttt{-len(}\textit{list}\texttt{)} to
    \texttt{len(}\textit{list}\texttt{)} \texttt{- 1}, inclusive, with negative
    indices wrapping around as in Python: if \textit{list}\texttt{[}$i$\texttt{]}
    is defined, then \textit{list}\texttt{[}$i$~\texttt{\%}~\texttt{len(}\textit{list}\texttt{)]}
    gives the same value.

    The operator \texttt{\textbf{*}} between a \texttt{list} and an \texttt{int}
    constructs a new list by repeating the elements of the given list the given
    number of times, with 0 or a negative number of times yielding
    \texttt{\textbf{[]}}; as opposed to Python, \texttt{\textbf{True}} and
    \texttt{\textbf{False}} are not allowed as the \texttt{int} operand. As in
    Python, the repetition makes shallow copies: \texttt{x = [[1,2]] * 4} makes
    \texttt{x[0] is x[1]} evaluate to \texttt{\textbf{True}}.

